\section{Experimental Setup}
\label{sec:setup}
\label{sec:experimental-setup}

\subsection{Workloads and Worker Traces}

We use a controlled sparse ridge-regression workload for the main experiments.
The generator creates a sparse feature matrix, partitions it into shards, and
computes one additive gradient contribution per shard.  A successful coded
decode therefore recovers the same full gradient as the uncoded shard sum, which
matches the linear aggregation model used by gradient coding~\cite{tandon2017gradient}
and sparse flexible coded learning~\cite{zhang2025sparseflexible}.  We use this
workload to isolate row-span recoverability, sparse row cost, and first-decode
timing, not to claim a full training-stack speedup for all sparse models.

Worker-state streams are controlled stress traces generated by the runtime.
Each iteration supplies per-worker speeds, straggler delays, and a slow-worker
mask.  We use four regimes: stable lognormal heterogeneity, bursty stragglers,
drifting base speeds, and phase changes in straggler fraction and slowdown.
These regimes model late-task effects in cluster systems~\cite{zaharia2008late}
and heterogeneous ML scheduling effects~\cite{narayanan2020gavel}.  The main
SoCC-facing experiments use phase-changing streams because the relative
importance of speed, delay, and row cost changes over time.  The traces are
controlled rather than production logs, so every paired comparison reuses the
same worker-state stream across strategies within a run.

\subsection{Metrics and Accounting}

Recent cluster schedulers report completion time, waiting time or slowdown,
resource utilization, and scheduling overhead~\cite{dongare2025rltune};
distributed-training planners additionally report throughput, cost, and search
time~\cite{strati2025sailor}.  Our setting has one synchronous coded-learning
iteration rather than a queue of multi-job training requests, so we adapt these
metrics to worker-service rounds.  The names below are the names used in the
evaluation tables and figures.
\begin{enumerate}
\setlength{\itemsep}{2pt}
\item \textbf{Mean barrier latency (\emph{Mean (ms)} or \emph{Barrier (ms)}).}
This is the end-to-end synchronous iteration time.  It includes first decode,
stream cancellation, and worker cleanup, which makes it the main metric for
deployed worker-service runs.  Stress tables use \emph{Barrier (ms)} for the
same quantity.
\item \textbf{Tail latency (\emph{p95 (ms)}).}  Tail latency is a standard
systems metric for cloud services~\cite{dean2013tail}.  We report p95 because
first-decode savings are useful only if they survive tail workers,
cancellation, and communication cleanup.
\item \textbf{Throughput (\emph{Thr. (iter/s)}).}  Training planners and cluster schedulers
commonly report throughput or goodput~\cite{strati2025sailor,qiao2021pollux}.
We use the reciprocal of barrier latency, reported as completed coded-learning
iterations per second.
\item \textbf{Worker-service overhead (\emph{Post rows} and \emph{Cancel (ms)}).}
Recent schedulers report scheduling runtime, search time, or planner overhead
alongside end-to-end performance~\cite{dongare2025rltune,strati2025sailor}.  In
our tables, \emph{Post rows} is the number of row results completed after the
first decodable event but before the round finishes, and \emph{Cancel (ms)} is
the cancellation/cleanup time included in the barrier path.  Lower values mean
less wasted work or less cleanup delay.
\item \textbf{Correctness and recovery counters (\emph{Succ.}, \emph{Gain},
\emph{Err.}, and \emph{Recov.}).}  Coded learning requires recovering the
intended gradient or row span despite
stragglers~\cite{tandon2017gradient,zhang2025sparseflexible}, while distributed
training systems also track worker faults and failure reactions~\cite{deng2025minder}.
\emph{Succ.} is the decode-success rate.  \emph{Gain} is the percentage
improvement over the comparison reference stated by the table or figure, such
as SFCL in the stress table.  \emph{Err.} is closed-worker errors per round, and
\emph{Recov.} is bounded row reissues per round.
\item \textbf{Mechanism and guard diagnostics.}  The row-exposure diagnostic
reports \emph{Mean gain}, \emph{p95 gain}, and $\Delta$ \emph{prefix}, all
relative to static placement in the same diagnostic.  Guard figures and tables
report \emph{decision rate}, \emph{false-enable rate}, \emph{false-disable
rate}, \emph{Mean/p95 $\tau$ error}, and \emph{Prefix-length error}.  Decision
rate is the fraction of rounds in which the guard enables adaptive scheduling;
false-enable and false-disable rates come from paired replay on the same
worker-state stream.
\end{enumerate}

These metrics test whether first-decode prefix savings survive scheduling,
communication, and cleanup costs.  Operationally, the scheduler should be
enabled only when the expected prefix saving exceeds the overhead budget,
\[
  \Delta T_{\mathrm{prefix}} >
  T_{\mathrm{sched}} + T_{\mathrm{extra}} + T_{\mathrm{cancel/comm}} .
\]

Replication counts, run identifiers, and raw per-run latency traces are
artifact details.  The main text keeps aggregate barrier, tail, throughput,
overhead, decision, and recovery results; the artifact retains the complete run
matrix, arm selections, paired comparisons, resource counters, and scripts that
regenerate the tables.

\subsection{Baselines and External Adapters}

The main comparison uses same-runtime baselines so that every method shares the
same sparse workload, worker-state stream, serialization path, cancellation
logic, and barrier accounting.  We use the following comparison points.
\begin{enumerate}
\setlength{\itemsep}{2pt}
\item \textbf{SFCL.}  This is the closest coded-computing baseline: static
sparse-flexible coded placement from SFCL~\cite{zhang2025sparseflexible}.  We
reimplement it in the same runtime.  It fixes the encoded-row assignment before
each round and uses neither worker heterogeneity nor first-decode prefix
features.
\item \textbf{RLTune-style adapter.}  This adapter follows the systems role of
RLTune~\cite{dongare2025rltune} by selecting among uncoded, speculative,
static-coded, and rank-aware coded arms using runtime heterogeneity features.
It is a same-runtime adapter, not a port of the full RLTune scheduler.
\item \textbf{Sailor-style adapter.}  This adapter follows the speed/cost
modeling role of Sailor~\cite{strati2025sailor}.  It assigns work using
worker-speed and row-cost estimates, but it does not use coded row-span or
first-decodable-prefix features.
\end{enumerate}

\textbf{Guarded enablement.}
The guard used in the evaluation is a deployment rule over runtime counters,
not a selector trained on final latency.  The online runtime estimates
decode-speed mismatch from row priorities and worker-speed estimates before an
iteration starts, predicts if the adaptive policy will shorten the
first-decodable prefix, and updates counters after the observed first-decode
event.  If the predicted prefix is worse, or recent counters show that adaptive
assignment expanded the completed-row prefix, \textsc{G-port} rejects the
adaptive arm and falls back to static sparse-flexible placement.  The portfolio
candidates share the same worker-service accounting; the ablations then remove
cross-policy execution or change the fallback rule.  The \texttt{best\_safe}
ablation chooses the cheaper predicted safe baseline between speed-aware
uncoded and static coded placement before the round.  Paired safe-baseline
replays over logged diagnostics estimate fallback headroom; they do not relabel
deployed K3s traces.  Unless stated otherwise,
the guard thresholds are $\theta_{\mathrm{cv}}=0.20$, $\theta_g=0.01$,
$\theta_K=0$ rows, and $\theta_a=1$ row; the EMA update weight is
$\alpha_{\mathrm{ema}}=0.30$ and the fallback decay is $\lambda=0.80$.

\subsection{Runtime Prototypes}

The experiments use three Python runtimes for sparse coded learning.  The
labels in Table~\ref{tab:external-main} are configurations of the
worker-service runtime: TCP runs local worker services over sockets,
TCP+stress adds analytic request/response delay and bandwidth limits to the
same TCP path, and K3s runs the same worker entrypoint as pods behind headless
Services.  The simulator and multi-process runtimes support mechanism checks;
the main external comparison uses the TCP-family and K3s worker-service
realizations.
\begin{enumerate}
\setlength{\itemsep}{2pt}
\item \textbf{Event simulator.}\par
The simulator generates sparse ridge data,
partitions it into shards, constructs the fixed sparse-flexible code matrix, and
samples worker completion times from the configured trace.  It runs the exact
row-span decoder after each row arrival.  We use it for mechanism sweeps,
oracle diagnostics, threshold sensitivity, and paired replays because it
separates scheduling logic from process, socket, and container overheads.

\item \textbf{Multi-process runtime.}\par
The local runtime starts one Python
worker process per logical worker.  The master assigns encoded row pairs,
workers compute sparse encoded gradients over their assigned shard
combinations, and completed rows are streamed back to the master.  After the
first decodable set is reached, the master cancels remaining work and records
both first-decode latency and conservative barrier latency, including
cancellation cleanup.  This runtime validates the master/worker control path
without network isolation.

\item \textbf{TCP worker-service runtime.}\par
The worker-service runtime launches
each worker as an independent TCP service.  Master and workers exchange
length-prefixed task, row, cancel, stop, and acknowledgement messages.  A task
message carries the round id, row ids, encoded rows, worker speed/delay state,
and current model vector; row messages return encoded gradients, row cost,
payload size, and timing counters.  The same worker entrypoint runs as local
services, host-published Docker containers, Docker-bridge containers, and
three-node K3s pods reached through headless Services.  This validates the
worker-service boundary without publishing worker ports.  For restricted cloud
hosts, the runtime also supports an SSH-forwarded two-node mode.  The K3s run is
a shared-pod validation, not an isolated production cluster; pod placement,
resource knobs, and manifests are artifact metadata.  The runtime also includes
the closed-connection stress hook: if a worker disconnects before first decode,
the master marks it unavailable, filters unfinished rows that have not arrived,
and reissues them to the fastest live worker.  Exact decoding remains the
correctness test; if reissued rows do not arrive in time, the iteration fails
normally and the next round can fall back to SFCL.
\end{enumerate}

The current runtimes are still prototypes: they de-risk the service boundary,
cancellation path, and small pod-to-pod TCP deployment, while production-scale
cluster scheduling and interference remain future systems work.  This is a
deployment step rather than a different algorithmic path, because the scheduler
only depends on streamed completions and cancellation acknowledgments.
